\subsection{Порядок выполнения}
\textbf{Первый проход:} прочитайте таблицу, назовите её служебные столбцы, затем выполните исходное сравнение без изменения настроек. Найдите одну строку результата и объясните её метрику. После этого переходите к исследованию гиперпараметров.

\subsubsection*{Шаг 1. Проверьте таблицы и метки}
Откройте файлы из \file{data}. Выведите размер таблиц, имена столбцов, число участников, распределение классов и наличие пропусков. Назовите два осмысленных признака. Не выбирайте входы срезом \file{values[:, :-2]}: положение служебных столбцов может измениться.
\begin{lstlisting}
from pathlib import Path
import pandas as pd
train = pd.read_csv(Path("data/train.csv"))
test = pd.read_csv(Path("data/test.csv"))
print(train.shape, test.shape)
print(train["Activity"].value_counts())
print(train["subject"].nunique(), test["subject"].nunique())
\end{lstlisting}

\subsubsection*{Шаг 2. Выделите валидацию}
Учебный модуль проверяет метки и признаки, исключает целевой столбец и идентификатор участника, затем выполняет групповое разделение. Если идентификатор называется иначе, укажите его через \file{group\_column}; не создавайте фиктивные номера участников из номеров строк.
\begin{lstlisting}
from lab1_utils import load_har, candidate_models, fit_validation
from lab1_utils import CLASSES, final_test, save_har_result
data = load_har("data", group_column="subject", seed=40)
print(data["Xtr"].shape, data["Xva"].shape)
print(data["training_subjects"], data["validation_subjects"])
\end{lstlisting}

\begin{keybox}{Контрольная точка}
Во входах нет \file{Activity} и идентификатора участника. Обучение, валидация и внешний тест содержат все шесть классов и разные группы людей. Валидация не участвует в вычислении параметров scaler. Столбцы теста приведены к порядку обучающих признаков по именам.
\end{keybox}

\subsubsection*{Шаг 3. Объясните выбор и выполните сравнение}
Сначала сформулируйте ожидаемое влияние гиперпараметров. В предлагаемом минимуме девять конфигураций с учётом константной модели. Допускается иной набор нескольких алгоритмов с сопоставимым и заранее определённым объёмом исследования.
\begin{table}[H]\centering\small
\caption{Исходный план сравнения классификаторов}
\begin{tabularx}{\textwidth}{@{}l l Y@{}}
\toprule
Обозначения & Модель & Меняемый параметр \\
\midrule
D0 & Dummy & Самый частый класс; ориентир без использования признаков \\
K3, K7, K15 & KNN & \(k\in\{3,7,15\}\); одинаковое масштабирование \\
S0.1, S1, S10 & LinearSVC & \(C\in\{0{,}1,1,10\}\); остальные настройки фиксированы \\
R8, RNone & Random forest & Глубина 8 или без ограничения; 200 деревьев \\
\bottomrule
\end{tabularx}\end{table}
\begin{lstlisting}
models = candidate_models(seed=40)
rows = []
for name, estimator in models.items():
    result = fit_validation(data, estimator)
    rows.append({"model": name, **result["metrics"],
                 "seconds_fit": result["seconds_fit"]})
comparison = pd.DataFrame(rows)
print(comparison.sort_values("macro_f1", ascending=False))
\end{lstlisting}

Этот цикл использует только обучение и валидацию. Сохраняйте предупреждения о сходимости; не скрывайте их общим отключением warnings. Если модель не сошлась, выясните причину и отдельно зафиксируйте изменённый бюджет или настройку.

\subsubsection*{Шаг 4. Визуализируйте и зафиксируйте выбор}
Постройте столбчатую диаграмму macro-F1 для всех конфигураций и отдельную диаграмму времени обучения. Не изображайте секунды и F1 как величины одной шкалы. Назовите худшие классы и предполагаемые причины ошибок. При близких результатах учитывайте сложность и время; не объявляйте небольшое различие устойчивым без повторной проверки.
\begin{lstlisting}
import matplotlib.pyplot as plt
comparison.plot.bar(x="model", y="macro_f1", legend=False)
plt.ylabel("Validation macro-F1")
plt.tight_layout()
\end{lstlisting}

\subsubsection*{Шаг 5. Выполните итоговый тест}
После выбора конфигурации переобучите её на всём \file{train.csv}; scaler внутри Pipeline также обучается заново. Только после этого получите прогноз \file{yhat} для \file{test.csv}. Пример ниже выбирает максимум macro-F1 валидации; окончательное объяснение выбора должно присутствовать в отчёте.
\begin{lstlisting}
selected = comparison.sort_values(
    ["macro_f1", "model"], ascending=[False, True]
).iloc[0]["model"]
final = final_test(data, models[selected])
yhat = final["prediction"]
print(pd.DataFrame(final["report"]).T)
save_har_result(data, comparison, selected, final, "runs/lab1_01")
\end{lstlisting}
Если результат теста оказался хуже ожиданий, опишите это как результат. Возвращение к выбору настроек после просмотра теста превращает его в ещё одну валидацию. Исходное требование сравнить несколько моделей выполняется на валидации; внешний тест оценивает уже выбранное решение.

\subsubsection*{Шаг 6. Разберите ошибки}
Постройте матрицу ошибок с одинаковым порядком шести меток. Выведите абсолютные числа и вариант с нормированием по строкам. Объясните хотя бы две заметные пары перепутанных действий; разделяйте наблюдение и гипотезу о его причине.
\begin{lstlisting}
from sklearn.metrics import ConfusionMatrixDisplay
ConfusionMatrixDisplay.from_predictions(
    data["ytest"], yhat, labels=range(6),
    display_labels=CLASSES, normalize="true", xticks_rotation=90
)
plt.tight_layout()
\end{lstlisting}

\subsection{Отчёт и вопросы защиты}
В отчёт включите паспорт файлов, состав разбиения по участникам, мотивацию выбора алгоритмов, таблицу всех конфигураций, две диаграммы, итоговый classification report и матрицу ошибок. Сохраните точные параметры модели, версии, seed и идентификаторы групп. Завершите выводом: какой вариант выбран, какие классы распознаются хуже и что можно проверить следующим опытом.

\begin{enumerate}
\item Что является объектом, признаком, целевой меткой и группой в этой работе?
\item Почему код активности нельзя трактовать как непрерывную величину?
\item Какую информацию дают акселерометр и гироскоп? Зачем нужны признаки частотной области?
\item Почему разбиение строк и разбиение участников проверяют разные условия применения?
\item Чем precision отличается от recall? Постройте собственный пример с FP и FN.
\item Как рассчитывается F1? Чем macro-F1 отличается от weighted-F1 и accuracy?
\item Чем гиперпараметр отличается от параметра обученной модели?
\item Как изменение \(k\), \(C\) и глубины деревьев повлияло на ваши результаты?
\item Почему scaler обучается только на обучающей части? Для каких моделей он особенно существенен?
\item Что означает нулевой support? Почему следует сохранять порядок меток в отчёте?
\item Почему нельзя выбрать победителя по внешнему тесту и затем назвать этот тест независимым?
\end{enumerate}

\begin{keybox}{Перед сдачей работы №\,1}
Есть сравнение нескольких моделей с вариациями гиперпараметров, визуализация, финальные предсказания и объяснение precision/recall/F1. Проверка отделена от подбора; \file{subject} не используется как входной признак. Указаны ограничения проверки на новых людях.
\end{keybox}
